\section{Introduction}

A final answer compresses a depth-wise process into one surface choice. The older draft in this repository tried to describe that process by comparing a weak first-hit event to stronger forever-after events. That comparison was asymmetric, so part of the reported lag could come from the metric definitions rather than from the score dynamics themselves.

This paper replaces that center with a fully persistence-based one. Under exact-option continuation scoring, we ask when the eventual winner first becomes a contender that never leaves the top-2 again, when it becomes the top-1 option for the rest of depth, and when it finally acquires a winner-runner gap that never again drops below half of its final value. Figure~\ref{fig:pooled-cdfs} shows those three persistent depths in one view.

The claim is deliberately narrow. We study only cached exact-option score trajectories for nine tested instruction-tuned open-weight models. We do not use this analysis to claim answer-format invariance, a hidden-state mechanism, or universality beyond the tested panel. Within that scope, the central pattern is consistent: the lag from persistent contender status to persistent decisive margin is typically nontrivial.
